\subsection{Reusing a head across generators}\label{sec:cross-head}

The head's inputs have the same geometric meaning in FM and GAGA. We test whether
its weights can also be shared: within each of the five fit pairs, we exchange
the trained heads verbatim and keep strength four. No target-generator force
labels train the transferred head, and no retraining follows the swap.
The comparison uses the same 64 compositions and initial draws as above,
without the optional H readout.

\begin{table}[t]\centering\small

\caption{Joint yield (percent) with the same or the other generator's trained head. Each cell contains 5,120 attempts. The transferred head uses only its source generator's force-training records.}

\begin{tabular}{lrrr}\toprule Target & No head & Own head & Other head\\\midrule
FM & 7.21 & 24.00 & 25.59\\
GAGA & 8.75 & 25.37 & 25.27\\
\bottomrule\end{tabular}\label{tab:cross-head}\end{table}

In the three new fit pairs, the FM-trained head improves GAGA joint yield by
16.54 percentage points (95\% composition interval $[13.02,20.08]$;
crossed fit/composition interval $[12.66,20.67]$). The three gains are
16.02, 18.36, and 15.23 points. Reverse transfer, from GAGA to FM, improves
joint yield by 17.12 points ($[13.61,20.77]$), with all three gains positive.
Thus the learned correction remains useful when the source distribution,
training objective, and sampling dynamics change. The experiment holds the
EGNN architecture and element vocabulary fixed.
Appendix~\ref{app:cross-head} gives every fit and the comparison to target-trained heads.
